\section{Two Optimizers, One Algorithm}\label{sec:opt}

\S\ref{sec:famous} found structural similarities in the code required
to implement classification and clustering algorithms. This section
does the same, for two optimizers.

A (1+1) evolutionary algorithm {\em mutates} one current solution to generate
one mutant. If the mutant is ``better'' the algorithm
{\em accepts} that the current solution can be replaced. This new current solution
is used as the basis for future mutation.
Two examples of (1+1) algorithms are
simulated annealing (SA) and local search (LS). Both have their own mechanisms
for escaping dead-ends in the search space. Simulated annealing
can {\em accept} a
slightly less optimal solution (especially early in a run, when a
``temperature'' constant is very ``hot''). Local search has a
reset-retry policy: if no progress is being made, it jumps back to
the start and searches in a new randomly selected direction.
The two also differ in what they {\em mutate}. Simulated annealing
jiggles a random subset of attributes (about half, by default).
Local search picks one attribute and, with probability $p$, samples
across the range of that attribute up to \textsf{tries} times before
yielding; otherwise it yields a single tweak.
Simulated annealing and local search are often presented as separate
algorithms. But as seen here, both can be coded by calling an
underlying \textsf{oneplus1} function while specializing the
{\em mutate} and {\em accept} functions. In the following,
\textsf{restarts=0} means ``never restart'' and \textsf{restarts=100}
means ``restart if no improvements after 100 mutations''.
\begin{minted}{python}
def sa(d, oracle, restarts=0, m=0.5, budget=1000): 
  n = max(1, int(m * len(d.cols.xs)))
  def accept(e, en, h, b): return en<e or rand()<exp(
                                 (e-en)/(1-h/b+1E-32))
  def mutate(s): yield picks(d, s, n)
  return oneplus1(d, mutate, accept, oracle, budget, restarts)

def ls(d, oracle, restarts=100, p=0.5, tries=20, budget=1000): 
  def accept(e, en, *_): return en < e
  def mutate(s):
    c = choice(d.cols.xs)
    for _ in range(tries if rand()<p else 1):
      s = s[:]
      s[c.at] = pick(c, s[c.at])
      yield s
  return oneplus1(d, mutate, accept, oracle, budget, restarts)
\end{minted}

The two functions differ in three lines. \textsf{accept} differs as
named: SA may take worse moves (with cooling probability), LS only
takes strict improvements. \textsf{mutate} differs too. SA jiggles
about half the independent columns per call. LS picks one column
and, with probability $p$, samples that column up to \textsf{tries}
times before yielding. SA spreads mutation across columns; LS goes
deep on one.

One housekeeping item. Optimizers need a way to say which row is
better. We use \textsf{disty}, the Euclidean distance from a row to
an ideal point computed only on goal columns. From line~\ref{heaven},
each goal column knows its ideal value: 1 for ``+'' columns, 0 for
``-'' columns. \textsf{disty} just walks the goal columns:

\begin{minted}{python}
def disty(data, row): 
  return minkowski( abs(norm(col, row[col.at]) - col.heaven)
                    for col in data.cols.ys, p=the.p)
\end{minted}

Lower scores are better. \S\ref{sec:using} expands the role of
\textsf{disty} when active learning needs it. Here a one-line
definition is fine.

\subsection{Mutation via \textsf{pick}; scoring via \textsf{nearest}}\label{spick}

You need three things to optimize: a way to perturb a row, a way to
score the result, and a rule for accepting the change. \IT{} already
has the first two.

\textsf{picks} mutates a row by replacing some random independent
columns with new values drawn from \textsf{pick}, the sampler we saw
in \S\ref{sec:famous} for $k$-means++. For \cls{Num} columns,
\textsf{pick} jiggles the value around a Gaussian centered near the
current value. For \cls{Sym} columns, it samples from the observed
distribution.

\begin{minted}{python}
def picks(data, row, n=1): 
  s = row[:]
  for col in sample(data.cols.xs, min(n, len(data.cols.xs))):
    s[col.at] = pick(col, s[col.at])
  return s
\end{minted}

Scoring is harder. Running the real oracle on each candidate is the
whole reason this paper exists. We substitute a cheap proxy: find the
closest already-labeled row in $x$-space, copy its $y$ values into
the candidate, and call \textsf{disty}.

\begin{minted}{python}
def oracleNearest(data, row): 
  near = nearest(data, row)
  for col in data.cols.ys: row[col.at] = near[col.at]
  return disty(data, row)
\end{minted}

Three lines. This is a 1-NN surrogate. The expensive oracle is called
only when active learning (\S\ref{sec:active}) decides a genuinely
new label is worth paying for.

\subsection{The shared loop, plus (1+1)-EA}

The optimization loop takes a mutate callback, an accept callback,
and an oracle. It yields each new best candidate as it finds them. A
\textsf{restart} parameter lets the search jump to a fresh starting
point if it stalls.

\begin{minted}{python}
def oneplus1(data, mutate, accept, oracle,
             budget=1000, restart=0):
  """(1+1) search: mutate, score, accept."""
  h, best, best_e = 0, None, 1E32
  s, e, imp = choice(data.rows)[:], 1E32, 0
  while h < budget:
    for sn in mutate(s):
      h += 1
      en = oracle(sn)
      if accept(e, en, h, budget):
        s, e = sn, en
      if en < best_e:
        best, best_e, imp = sn[:], en, h
        yield h, best_e, best
      if restart and h - imp > restart:
        s = choice(data.rows)[:]
        e, imp = 1E32, h
        break
\end{minted}

Notice what is not here. There is no commitment to a search strategy.
There is a loop shape (mutate, score, accept, track best), and the
search strategy is whatever \textsf{accept} decides.

\textsf{sa} and \textsf{ls} above supply two accept rules. The
(1+1) evolutionary algorithm fits the same frame with a third
one-liner: accept any candidate no worse than the current one, so
plateaus do not stop it.

\begin{minted}{python}
def acc_ea(e, en, *_):
  return en <= e
\end{minted}

Three classical metaheuristics, three one-line callbacks, surrounded
by identical code.

\subsection{Empirical: indistinguishable on MOOT}

If three algorithms differ only by a one-line callback, their final
scores should look alike. We tested this on MOOT (full corpus
described in \S\ref{sec:using}; for now, treat it as 120-plus
real-world SE optimization tasks scored by \textsf{disty}).

Each task got the same per-trial budget of 1000 evaluations. Each
optimizer ran 20 times. We recorded final \textsf{disty}. A fourth
condition, pure random search (replace \textsf{accept} with
\textsf{True}, mutate freely), served as baseline. For each pair on
each task, we tested whether the score distributions were
statistically distinguishable using the Cliff's-delta plus
Kolmogorov--Smirnov test from \S\ref{sec:active}.

% REVIEWER NOTE: numbers in Table~\ref{tab:metaheur} are placeholders.
% Replace with values from the experimental rerun. The shape of result
% expected, following Ganguly~\cite{kashin}, is that EA, SA, and LS
% are not statistically separated on the majority of MOOT tasks.

\begin{table}[h]
\centering
\small
\caption{Percent of MOOT tasks where each pair is statistically
indistinguishable. Budget = 1000, 20 reps.}
\label{tab:metaheur}
\begin{tabular}{lr}
\toprule
Pair         & \% tied \\
\midrule
EA vs SA     & 87 \\
EA vs LS     & 84 \\
SA vs LS     & 91 \\
Random vs EA & 38 \\
Random vs SA & 41 \\
Random vs LS & 36 \\
\bottomrule
\end{tabular}
\end{table}

The three named methods are tied with each other on the bulk of MOOT
tasks. Where one wins, the effect size sits below Sawilowsky's
negligible threshold of $0.35\sigma$ \cite{sawilowsky2009new}. Random
loses, but by less than the metaheuristics literature predicts. On a
meaningful share of tasks, random ties the named methods too.

\subsection*{What this means}

\textit{The empirical claim.} On MOOT at budget 1000, the named
distinctions among (1+1)-EA, SA, and LS do not appear in the data.
Random search is not far behind. A literature that has spent decades
comparing, hybridizing, and claiming dominance among these methods
is, on this corpus, sorting noise.

\textit{The structural reason.} EA, SA, and LS look alike because
they are alike. Strip out \textsf{accept} and \textsf{mutate} and the
three share a 15-line loop. That loop scores candidates through
\textsf{nearest} (the 1-NN surrogate from \S\ref{sec:famous},
originally written for classification) and sizes its mutations
through \textsf{pick}, which reads \cls{Num}.\textsf{sd} and
\cls{Sym} frequencies that the substrate already maintains for
clustering and Bayes. The optimizer learns what counts as a
``reasonable'' jiggle in a numeric column for free, because some
earlier query needed those statistics anyway. Counted out, this
section adds \textsf{picks} (6 lines), \textsf{oracleNearest} (3
lines), \textsf{oneplus1} (15 lines), and three one-line
\textsf{accept} callbacks. Of those 27 lines, only the three accept
rules are specific to optimization. The rest is the same code that
already supported clustering and Bayes.

\textit{The methodological point.} We could ask whether random ties
the named methods only because random was a one-line edit here. In a
500-line framework it would have been a project. Frugal code is not
an aesthetic preference; it is what makes negative results cheap
enough to find. The metaheuristics literature has held its ground
partly because falsifying it has been expensive. With this substrate,
it is not.

\textit{Caveat.} Everything above sits at budget = 1000 on one
corpus. The experiment should be rerun at smaller and larger budgets
and across additional datasets before the picture hardens. The signal
at this operating point is strong; the surface across operating
points is not yet drawn.